\section{Local rigidity and the corrected boundary theorem}\label{sec:cold}

For a block $P\subseteq[X,2X]$, write $N=|P|$ and
\[
 Q_P(a)=\sum_{p<q\in P}\left(\frac{H_{pq}(a)}{pq}\right)^2,
 \qquad
 \sigma_P^2=\sum_{p<q\in P}\frac1{(pq)^2}.
\]
Fix a dominance parameter
\begin{equation}\label{eq:rho-range}
 0<\rho<\frac12.
\end{equation}
An integer $m$ is a dominant label if it lies in the stated centred size range
and
\begin{equation}\label{eq:dominant}
 \#\{p\in P:a_p\equiv m\pmod p\}\ge(1-\rho)N.
\end{equation}
Two dominant labels are equal once the block is large enough: their agreement
classes contain two common primes, so the product of those primes divides the
difference of the labels, whose centred size is smaller than that product.

\begin{proposition}[Nondominant forcing]\label{prop:nondominant}
For every $\rho$ satisfying \eqref{eq:rho-range}, there are constants
$c_w(\rho)>0$ and $X_0(\rho)$ such that, whenever $X\ge X_0(\rho)$,
$P\subseteq[X,2X]$ is a prime set with $N\ge X/(2\log X)$, and no dominant label
exists, then
\begin{equation}\label{eq:nondominant-floor}
 Q_P(a)\ge c_w(\rho)\frac{X}{\log^3X}.
\end{equation}
\end{proposition}
\begin{proof}
We record the covering-and-dispersion mechanism because it is used again in the
global count.  If $Q_P(a)\le R$, choose a threshold
\[
 B\asymp_\rho \sqrt R\,\frac{X^2}{N}.
\]
Only $O(RX^4/B^2)$ pairs can have $|H_{pq}|>B$.  Double counting therefore
produces a base prime $p_0$ with few such neighbours.  For every other
nonexceptional $q$, the integer $H_{p_0q}$ is a label of size at most $B$, has
the required residue modulo $q$, and belongs to one residue class modulo $p_0$.
There are consequently $O(1+B/X)$ possible labels.

Partition the nonexceptional primes by these labels.  If no class has size
$(1-\rho)N$, then either a positive proportion of the primes lies in substantial
classes outside a largest class, or a positive proportion lies in small
classes.  In the first case, cross-label dispersion and power mean give
$R\gg_\rho X/\log^3X$; in the second, the number and sizes of the small classes
give the same lower bound.  Increasing $X_0(\rho)$ secures the auxiliary
integer and size inequalities.  This proves \eqref{eq:nondominant-floor}.
\end{proof}

Call block $k$ \emph{hot} if
\begin{equation}\label{eq:forcing-floor}
 Q_k(a)\ge R_w(k):=c_w(\rho)\frac{2^k}{\log^3(2^k)},
\end{equation}
and \emph{cold} otherwise.  Every sufficiently large cold block has a unique
dominant label.  Let $C_m$ be its agreement class and $E=P\setminus C_m$ its
exception set.  The same dispersion argument, now with one fixed dominant
class, gives constants depending only on $\rho$ such that
\begin{equation}\label{eq:exception-count-symbolic}
 |E|\ll_\rho \frac{R\log^3X}{X}
\end{equation}
and
\begin{equation}\label{eq:label-bound}
 |m|\ll_\rho\frac{\sqrt R}{\sigma_P},
 \qquad
 \sigma_P\gg\frac{N}{X^2}.
\end{equation}
Choose the forcing normalization $c_w(\rho)$ sufficiently small, before the
bottom scale, so that the exception count lies within a fixed budget $e_0$ and
vanishes in the later localization regime.  At a common scale threshold the
label-size estimate also gives
\begin{equation}\label{eq:boundary-label-size}
 32|m|\le X(N-e_0),
 \qquad N-e_0\ge12,
\end{equation}
with the analogous inequalities in the next block.

The boundary theorem requires all of these hypotheses.  Distinct adjacent
integer labels alone do not force energy.

\begin{proposition}[Exception-aware adjacent-label penalty]
\label{prop:boundary-penalty}
Let consecutive cold blocks $P_k,P_{k+1}$ have dominant labels $m\ne m'$ and
exception sets $E_k,E_{k+1}$ of cardinality at most $e_0$.  Assume residue
agreement with the labels outside the named exception sets, the reduced
cardinality bounds, and the label-size bounds in
\eqref{eq:boundary-label-size} and their upper-block analogues.  Then the
bipartite energy $X_k(a)$ satisfies
\begin{equation}\label{eq:Pi}
 X_k(a)\ge\Pi_k:=
 \frac{(|P_{k+1}|-e_0-1)(|P_k|-e_0)^3}
      {2^{13}(2^k)^2}.
\end{equation}
\end{proposition}
\begin{proof}
Put $d=m'-m$.  For $p\in P_k$, $q\in P_{k+1}$, and an inverse $\bar p$ modulo
$q$, the centred CRT representative satisfies
\[
 \normdist{d\bar p/q}=\normdist{(H_{pq}-m)/(pq)}.
\]
The nearest-integer triangle inequality and the label-size bound absorb the
$m$-term.  For every nonexceptional upper prime $q\nmid d$, cross-block
dispersion gives
\[
 \sum_{p\in P_k\setminus E_k}
 \left(\frac{H_{pq}}{pq}\right)^2
 \ge\frac{|P_k\setminus E_k|^3}{2^{13}(2^k)^2}.
\]
At most one nonexceptional upper prime divides $d$: the product of two such
primes would exceed $|d|$ while dividing it.  Summing over the remaining upper
primes gives \eqref{eq:Pi}.
\end{proof}

The local powers of two and rational constants used in one implementation of
the covering, exception, and boundary estimates are finite certificates.  The
structural content retained here is the forcing scale, the quantitative
exception and label bounds, and every hypothesis of Proposition
\ref{prop:boundary-penalty}.  Without those hypotheses the assertion is false:
for example, a label equal to the product of all lower-block primes can represent
the same lower residues as zero.
